We consider a generalized plant $P$ with a PIE state-space realization.
The nominal PIE duality results of Shivakumar et al.~\cite{shivaknew}
provide the starting point for convex state-feedback synthesis.
For dynamically filtered IQCs, two further ingredients are relevant,
hard factorizations connecting IQC theory to
dissipativity theory~\cite{9781263,6915700,scherer2018dynamicdissipation}, and the dual system representations
developed for LPV estimation and feedforward
synthesis~\cite{kose2009robust,venkataraman}.
This section develops the corresponding dual-filter construction for PIE
systems using the induced $L_2$-gain PIE duality result.

\subsection{Nominal PIE induced $L_2$-norm duality}\label{sec:dynamic-filtering-duality-1}
We begin by restating the main theorem from \cite{shivaknew}. In \cite{shivaknew}, 
the theorem is formulated for the case where $\mcl{D}$ is a matrix and the 
induced $L_2$ norm is bounded over an infinite time horizon. Generalizing this result 
to the case where $\mcl{D}$ is a PI operator and considering the bounded truncated 
induced $L_{[a,b]}$ norm does not alter the proof.


\begin{theorem}[Dual $L_2$-gain
{\cite[Thm.~9]{shivaknew}}]\label{thm:dual-l2-gain}
For a PIE system with corresponding PI operators $P\{\mcl{T}, \mcl{A}, \mcl{B},
\mcl{C}, \mcl{D}\}$, the following statements are equivalent.
\begin{enumerate}
\item For any $w\in L_{\mrm{e},[a,b]}^{\mbf{n}_{\mrm{w}}}$, and
$\mbf{x}(0)=0$, $P$
satisfies $\|p_T {z}\| \leq \gamma \|p_T {w}\|$ for all $T>0$.
\item For any $\underline{w}\in L_{\mrm{e},[a,b]}^{\mbf{n}_{\mrm{z}}}$, and
$\underline{\mbf{x}}(0)=0$, $P^\top$ satisfies $\|p_T \underline{z}\| \leq
\gamma \|p_T \underline{w}\|$ for all $T>0$.
\end{enumerate}
\end{theorem}
For a bounded, time-invariant, memoryless linear uncertainty
$\Delta$, suppose the feedback interconnection $P\star\Delta$ is well-posed,
in accordance with~(\ref{ass:well-posedness}).
Closing the loop $w=\Delta z$ gives
\[
\mcl T\dot{\mbf x}
=\bigl(\mcl A+\mcl B\Delta(I-\mcl D\Delta)^{-1}\mcl C\bigr)\mbf x.
\]
Taking adjoints and using the push-through identity yields
\begin{equation*}
(\mcl{T})^*=\mcl{T}^*, \quad\begin{aligned}
&\bigl(\mcl A+\mcl B\Delta(I-\mcl D\Delta)^{-1}\mcl C\bigr)^*\\
&\quad=\mcl A^*+\mcl C^*\Delta^*(I-\mcl D^*\Delta^*)^{-1}\mcl B^*.
\end{aligned}
\end{equation*}
This is exactly the state operator of $P^\top\star\Delta^*$. Thus, \cite[Thm.~6]{shivaknew} establishes
equivalence of weak convergence for the primal and dual states. 
Exponential stability is likewise equivalent
by \cite[Cor.~8]{shivaknew}. For the constant repeated real uncertainty
$\Delta=\delta I$ in Section~\ref{ex:coupled-heat},
$\Delta^*=\Delta$, so the same uncertainty structure appears in both interconnections, provided
the performance IQC $(\Psi_\mrm{p}, \mcl{V}_\mrm{p})$ admits a J-spectral factorization.
\subsection{Duality under dynamic filtering}
For a general uncertainty $\Delta$, this construction is
less trivial. It is therefore more convenient to develop duality
through the IQC description, whose filters and multiplier weights are
linear operators, rather than through $\Delta$ itself.
The inertia argument of~\cite{kose2009robust} assumes linear uncertainty.
Venkataraman and Seiler~\cite{venkataraman} accommodate nonlinear
uncertainties through dual IQCs and $J$-spectral factorizations of strict
Positive--Negative (PN) IQCs. We adopt the same strategy.
We first reformulate the IQC condition as a strict small-gain requirement 
on an associated system \(G_0\) with respect to the 
\(L_{e,[a,b]}\) norm, then apply Thm.\ref{thm:dual-l2-gain},
and finally reconstruct the dual system. 
In their framework, this relationship is established through congruence 
transformations of the analysis LMIs~\cite[Sec.~2-3]{venkataraman}. Our approach relies 
on PIE induced $L_{e,[a,b]}$-gain duality result. Let us define a hard J-spectral factorization.
\begin{definition}[$J$-Spectral Factorization and Dual Filter] 
\label{def:dynamic-filters}\label{def:factorization}
For a given IQC defined by $(\Psi,{\mcl V})$, we say that  a causal linear bounded operator $\Theta$ and $J_{\mbf n_{\mrm{z}},\mbf n_{\mrm{w}}}:=\bmat{ I_{\mbf n_{\mrm{z}}} & 0 \\ 0& -I_{\mbf n_{\mrm{w}}} }$
(the pair $(\Theta,J_{\mbf n_{\mrm{z}},\mbf n_{\mrm{w}}})$) is a $J$-spectral factorization of $(\Psi,{\mcl V})$ if, 
\begin{equation*}
\ip{\Psi\bmat{z\\w}}
{\mcl M_{\mcl V}\Psi\bmat{z\\w}}_{L_{[a,b]}}>0
\end{equation*}
can be factorized as
\begin{equation*}
\ip{p_T\Theta\bmat{z\\w}}
{p_TJ_{\mbf n_{\mrm{z}},\mbf n_{\mrm{w}}}\Theta\bmat{z\\w}}_{L_{[a,b]}}>0 \quad \forall T>0,
\end{equation*}
Such that $\Theta,\Theta_{22}$ are causally boundedly invertible. Following the 
dual-filter construction in~\cite{venkataraman}, we use $\mbf D(\Theta)$ to 
denote the dual filter associated with $\Theta$ of the form
\begin{equation*}
\mbf{D}(\Theta)
:=
\left(\hat{J}^*\Theta^\top\hat{J}\right)^{-1},
\end{equation*}
where $\hat{J}:=\bmat{0_{\mbf n_{\mrm{z}}\times\mbf n_{\mrm{w}}}&I_{\mbf n_{\mrm{z}}}\\-I_{\mbf n_{\mrm{w}}}&0_{\mbf n_{\mrm{w}}\times\mbf n_{\mrm{z}}}}$.
\end{definition}

\begin{remark}[Hard factorization requirement]
For finite-dimensional channels, the hard-factorization requirement covers a
broad class of multipliers~\cite{9781263,6915700,scherer2018dynamicdissipation}. In particular, strict PN
multipliers, even when introduced through soft frequency-domain IQCs, admit a
hard $J$-spectral factorization~\cite[Thm.~4]{6915700}. Strict PN is also
preserved by the primal--dual transformation~\cite{venkataraman}.
For infinite-dimensional channels,
however, hardness of a proposed $J$-spectral factorization must be established
explicitly. Callegari et al.~\cite{callegari} propose IQC lifting, which could
extend finite-dimensional rational $J$-spectral factorizations to multipliers
acting on infinite-dimensional channels.
\end{remark}
For further brevity, define the primal quadratic form
\[\mcl{I} := \ip{p_T \Theta\bmat{Pw \\ w}}
{p_T J_{\mbf{n}_{\mrm{z}}, \mbf{n}_{\mrm{w}}}\Theta\bmat{Pw \\ w}}_{L_{[a,b]}} .\]
The corresponding dual form is evaluated on the graph of $P^\top$ using the
dual filter $\mbf D(\Theta)$ and the $J_{\mbf{n}_w, \mbf{n}_z}$ operator,
\[\underline{\mcl{I}} := \ip{p_T \mbf{D}(\Theta)\bmat{P^\top\underline{w} \\ \underline{w}}}
{p_T J_{\mbf{n}_{\mrm{w}}, \mbf{n}_{\mrm{z}}}\mbf{D}(\Theta)\bmat{P^\top\underline{w} \\ \underline{w}}}_{L_{[a,b]}} .\]
The following theorem is our \emph{main result} and establishes the system relations in Figure~\ref{fig:potapov-ginzburg-commutative-diagram} 
\begin{theorem}[Primal--dual IQC equivalence]\label{thm:dual-iqc}
Let $P$ be as in Definition~\ref{def:uncertain-interconnection}, and let
$(\Theta,J_{\mbf n_{\mrm z},\mbf n_{\mrm w}})$ be a hard
$J$-spectral factorization. Assume that both $P$ and $\Theta$ admit
PIE realizations.    Finally, assume $(\Theta_{21}P +\Theta_{22})$ has a bounded causal inverse. 
    Define 
    \begin{equation*}
        G_0 = (\Theta_{11}P +\Theta_{12})(\Theta_{21}P+\Theta_{22})^{-1}
    \end{equation*} 
    then the following statements are equivalent,
    \begin{enumerate}
        \item There exists some $\rho<1$ such that $\|p_T G_0 \tilde{w}\|^2_{L_{[a,b]}}\leq\rho\|p_T \tilde{w}\|^2_{L_{[a,b]}}$ for every $\tilde{w}$ and $T\geq 0 $.
        \item There exists some $\epsilon > 0$ such that $\mcl{I}\leq - \epsilon \|p_T w\|_{L_{[a,b]}}^2$ for every $w$ and $T\geq0$.
        \item There exists some $\underline{\epsilon}>0$ such that $\underline{\mcl{I}} \leq - \underline{\epsilon}\|p_T\underline{w}\|_{L_{[a,b]}}^2$ for every $\underline w$ and $T\geq0$.
    \end{enumerate}
\end{theorem}
\begin{proof}
    See \ref{proof:dual-iqc}.
\end{proof}
\subsection{Dual robust-performance certificate}

Theorem~\ref{thm:dual-iqc} establishes equivalence of the strict primal and dual IQC
inequalities through a small-gain requirement on $G_0$. Either side of the
filtered graph may therefore be used to construct the same performance
certificate. The resulting
KYP condition permits the robust-performance certificate to be imposed on
either the primal or dual realization. 
\begin{corollary}\label{thm:dual-KYP}
   Let $P$ be as in Definition~\ref{def:uncertain-interconnection}, and let
$(\Theta,J_{\mbf n_{\mrm z},\mbf n_{\mrm w}})$ be a hard
$J$-spectral factorization. Assume that both $P$ and $\Theta$ admit
PIE realizations. Let $\Delta \in \mbf{\Delta}$ such that, 
    \begin{itemize}
        \item The interconnection $P\star\Delta$ is well-posed for all
        $\Delta\in\mbf{\Delta}$.
        \item The IQC in Definition~\ref{def:uncertainty-iqc} is satisfied for all $\Delta\in\mbf{\Delta}$.
        \item Let $\mcl D$, $\mcl{D}_\mrm{P}$ denote the feedthrough of $\Theta$ and $P$ respectively, such that $(\mcl{D}_{21}\mcl{D}_\mrm{P} + \mcl{D}_{22})^{-1}\in\Pi_4$.
    \end{itemize}
    Suppose that \emph{either} of the two statements hold for some $\mcl{P} = \mcl{P}^* \succeq 0 \in \Pi_4$.
    \begin{itemize}
        \item \begin{minipage}[c]{\linewidth}
        \begin{equation*}
            \bmat{\mcl T^*\mcl P\mcl A + (*) & \mcl T^*\mcl P\mcl B\\
            \mcl B^*\mcl P\mcl T & \epsilon I} + \bmat{\mcl C^* \\ \mcl D^*}J_{\mbf{n}_z, \mbf{n}_w}\bmat{\mcl C & \mcl D} \preceq 0
        \end{equation*}
        \end{minipage}
        \item \begin{minipage}[c]{\linewidth}
        \begin{equation}\label{eq:dual-kyp}
            \bmat{\underline{\mcl T}^*\mcl P\underline{\mcl A} + (*)& \underline{\mcl T}^*\mcl P\underline{\mcl B}\\
            \underline{\mcl B}^*\mcl P\underline{\mcl T} & \underline{\epsilon} I} + \bmat{\underline{\mcl C}^* \\ \underline{\mcl D}^*}J_{\mbf{n}_w, \mbf{n}_z}\bmat{\underline{\mcl C} & \underline{\mcl D}} \preceq 0
        \end{equation}
        \end{minipage}
    \end{itemize}
    where, for any input $w\in L_{\mrm{e},[a,b]}$ or
    $\underline w\in L_{\mrm{e},[a,b]}$, the corresponding output
    $z$ or $\underline z$ satisfies the primal or dual realization,
    \begin{equation*}
        \Theta\bmat{P\\I}w := \left\{ \bmat{\mcl T \dot{\mbf{x}}(t) \\ v(t)} = \bmat{\mcl A & \mcl B\\
        \mcl C & \mcl D}\bmat{\mbf{x}(t) \\ w(t)}\right.
    \end{equation*}
    or
    \begin{equation}\label{eq:dual-system}
        \mbf{D}(\Theta)\bmat{P^\top\\I}\underline{w} := \left\{ \bmat{\underline{\mcl T} \dot{\underline{\mbf{x}}}(t) \\ {\underline{v}}(t)} = \bmat{\underline{\mcl A} & \underline{\mcl B}\\
        \underline{\mcl C} & \underline{\mcl D}}\bmat{\underline{\mbf{x}}(t) \\ \underline{w}(t)}\right.
    \end{equation}
    Then the interconnection $P\star\Delta$ with $\mbf{x}(0) = 0$, as
    in~(\ref{ass:initial-conditions}), is robustly performing.
\end{corollary}
\begin{proof}
    See \ref{proof:dual-KYP}
\end{proof}
The next section applies duality to recover convex state-feedback synthesis.
